% =====================================================================
% IEEE jurnal maqolasi shabloni (IEEEtran sinfi) - Ilmiy Yo'l
% Izohlar o'zbek tilida, maqola matni inglizcha yoziladi.
%
% Kompilyatsiya: pdflatex -> bibtex -> pdflatex -> pdflatex
% (Overleaf'da "Recompile" tugmasi buni o'zi bajaradi.)
%
% MUHIM: yuboriladigan jurnalning "Information for Authors" sahifasini
% o'qing. Ba'zi jurnallar [journal] o'rniga [conference] yoki boshqa
% sozlama talab qiladi.
% =====================================================================
\documentclass[journal]{IEEEtran}

\usepackage[T1]{fontenc}
\usepackage[utf8]{inputenc}
\usepackage{cite}
\usepackage{amsmath,amssymb}
\usepackage{graphicx}
\usepackage{booktabs}
\usepackage{url}

\hyphenation{op-tical net-works semi-conduc-tor}

\begin{document}

% --- Sarlavha: aniq, mazmunli, qisqa. Qisqartmalardan qoching. ---
\title{Title of Your Article: Specific, Informative and Concise}

% --- Mualliflar va ish joyi ---
\author{First~Author,~\IEEEmembership{Member,~IEEE,}
        Second~Author,
        and~Third~Author%
\thanks{Manuscript received Month DD, 2026; revised Month DD, 2026.}%
\thanks{First Author and Second Author are with the Department of Information Technologies, Example University, City, Country (e-mail: author@example.com).}%
\thanks{Third Author is with the Example Institute, City, Country.}}

\markboth{IEEE Journal Name, Vol.~XX, No.~X, Month 2026}%
{Author \MakeLowercase{\textit{et al.}}: Title of Your Article}

\maketitle

% --- Abstract: muammo, usul, asosiy natija, xulosa. Bitta xatboshi. ---
\begin{abstract}
State the problem in one sentence. Describe the method you propose or the data you analyse. Report the main quantitative result, for example the improvement over the baseline, with a number. End with the practical or scientific significance. Keep the abstract self-contained, without citations or abbreviations that are not explained.
\end{abstract}

\begin{IEEEkeywords}
keyword one, keyword two, keyword three, keyword four.
\end{IEEEkeywords}

\IEEEpeerreviewmaketitle

% =====================================================================
% I (Introduction): muammo nima, nima uchun muhim, boshqalar nima qilgan,
% sizning ishingiz nimasi bilan yangi (research gap), maqola tuzilishi.
% =====================================================================
\section{Introduction}
State the research problem and why it matters to the field~\cite{he2016resnet,redmon2016yolo}. Summarise what earlier studies have done and where they fall short. Name the gap clearly: new data, a new method, or a new context.

State the contribution of this article in one sentence, then list the contributions if there are several. Finish with a short outline of the remaining sections.

% =====================================================================
% M (Methods): boshqa olim ishingizni takrorlay oladigan darajada aniq.
% =====================================================================
\section{Methodology}
Describe the data set, its source, size and how it was collected or prepared. Describe the proposed method step by step. Give the evaluation metric in a formula, for example the $F_1$ score:
\begin{equation}
  F_1 = \frac{2PR}{P+R},
  \label{eq:f1}
\end{equation}
where $P$ is precision and $R$ is recall. Refer to equations in the text as~(\ref{eq:f1}).

% =====================================================================
% R (Results): faqat natijalar, izohsiz. Jadval va rasm matnda tilga olinadi.
% =====================================================================
\section{Results}
Present the results without interpretation. Table~\ref{tab:results} shows the comparison with the baseline. Replace the placeholders with your own measured values.

\begin{table}[!t]
  \caption{Comparison With the Baseline (Replace With Your Results)}
  \label{tab:results}
  \centering
  \begin{tabular}{lcc}
    \toprule
    Method & $F_1$ score (\%) & Time (ms) \\
    \midrule
    Baseline & XX.X & XX \\
    Proposed & XX.X & XX \\
    \bottomrule
  \end{tabular}
\end{table}

% Rasm qo'shish uchun (fayl papkada bo'lishi kerak, masalan fig1.pdf):
% \begin{figure}[!t]
%   \centering
%   \includegraphics[width=\columnwidth]{fig1.pdf}
%   \caption{Short, self-explanatory caption.}
%   \label{fig:overview}
% \end{figure}

% =====================================================================
% D (Discussion): natijalar nimani anglatadi, cheklovlar, boshqa ishlar bilan solishtirish.
% =====================================================================
\section{Discussion}
Interpret the results. Explain why the proposed method behaves as observed, compare with related work~\cite{vaswani2017attention} and state the limitations honestly.

\section{Conclusion}
Summarise the findings in two or three sentences and state the future work. Do not repeat the abstract word for word and do not introduce new results.

\section*{Acknowledgment}
Thank the people and funding sources that supported the work. Write ``None'' if there are none.

% --- Adabiyotlar: references.bib faylidan. Faqat matnda keltirilganlari chiqadi. ---
\bibliographystyle{IEEEtran}
\bibliography{references}

\end{document}
